\cotefile{88}
\scanpage{2}
\begin{flushright}
\unsure{Cartes|Carte} associées aux syst. de Racines

(7)
\end{flushright}
% archive sheet number: 1

\scanpage{3}
\diagram{\begin{tikzcd}[ampersand replacement=\&,column sep=large]
\text{Cat des systèmes de racines}\arrow[r,"\Phi"] \& \text{cat. des systèmes}
\end{tikzcd}}
\[
(W,R,\deleted{\unsure{\varphi}}I,\varphi:R\times I\to W)
\]
où $W$ est un groupe, $R$ un $W$-torseur à g., \inserted{$I$ ens. fini,} $\varphi:R\times I$ une \deleted{\unsure{de|$I$ $W$}} application

avec

a) $\varphi(wr,i)=\operatorname{int}(w)\varphi(r,\unsure{i'|i})$

b) $\forall r\in R,\quad i\mapsto\varphi(r,i)$ définit un système de générateurs de $W$ qui en fait \deleted{\unsure{de}} un \textquotedblleft groupe de Weyl épinglé\textquotedblright\ i.e.

c)
\[
\left\{\begin{aligned}\varphi(r,i)^2&=1\\(\varphi(r,i)\varphi(r,j))^{n_{ij}}&=1\quad(i\ne j)\end{aligned}\right.
\]
est une \uline{présentation} de $W$ (où $n_{ij}$ est l'ordre de $\varphi(r,i)\varphi(r,j)$), et

d) $I$ muni \deleted{\unsure{des}} \inserted{du} syst. des $n_{ij}$ ($n_{ij}=n_{ji}$) est un diagramme de Dynkin ($i$ lié à $j$ sss $n_{ij}\geq 3$ i.e. $n_{ij}\ne2$ i.e. $\varphi(r,i)$ ne commute pas à $\varphi(r,j)$)

\uline{NB} On a
\[
 R\xrightarrow{\psi}\ \text{Syst. gén admissibles de }W
 \quad\text{paramétrés par le diagramme de Dynkin }I
\]
Mais cette application, \deleted{\unsure{il se peut}} est \unsure{bien} \unsure{a)} compatible aux actions de $W$, n'\deleted{\unsure{est}}est pas \inserted{\unsure{tjrs|en gén.}} injective, \inserted{si $r\in R$, elle est inj.} \deleted{(si $I\ne\varnothing$), \unsure{car en fait} si $r,r'$ sont deux \textquotedblleft\unsure{repères}\textquotedblright\ correspondant à des syst. simples de racines opposés l'élément \inserted{$w$} du groupe de Weyl qui \unsure{amè|envoie}} ssi le stabilisateur de \unsure{$\psi$|$\psi(r)$|repères} ds $W$ est réduit à l'élément neutre. Or ce stabilisateur est évidemment le \uline{centre} de $W$ -- donc $\psi$ injective (permet d'interpréter $R$ comme un ens de syst. de gén. adm. de $W$ paramétrés par le diagramme de Dynkin $I$) \deleted{\illegible{1}} sss $Z(\unsure{W|w})=\{1\}$.

Le foncteur précédent est-il une équivalence de catégories ? Il est \inserted{\unsure{évid.\textsuperscript{t}|évid.|évid. fid.}} \unsure{ess\textsuperscript{t}|très|ess.} surjectif, est-il pl. fidèle ? Deux syst. de racines $\Sigma,\Sigma'$ tels que $\Phi(\Sigma)\sim\Phi(\Sigma')$ ont \deleted{\unsure{di}} gr. de Dynkin isom,
% archive sheet number: 2

\scanpage{4}
donc sont isomorphes. Il reste à voir que pour tt $\Sigma$,
\[
\operatorname{Aut}(\Sigma)\xrightarrow{\unsure{\prime|\sim}}\operatorname{Aut}(\Phi(\Sigma))
\]
est un iso de groupes. Or on a un hom. canonique
\[
W\xrightarrow{\alpha}\operatorname{Aut}(\Phi(\Sigma))
\]
par
\[
\alpha(w)=(\operatorname{int}(w),w_R,\operatorname{id}_I)
\]
[en effet,
\[
\begin{aligned}
w_R(\deleted{\unsure{w}}\,\underset{\substack{\uparrow\\W}}{w'}\underset{\substack{\uparrow\\R}}{r})&=\operatorname{int}(w)(w')\,w_R\,r\\
\Vert\qquad &\qquad\Vert\\
w(w'r)&\qquad(ww'w^{-1})wr
\end{aligned}
\]
et
\[
\left.\begin{aligned}
\varphi(w_R(r)\deleted{\unsure{r}},i)&=\operatorname{int}(w)\varphi(r,i)\\
\Vert\qquad &\\[-2pt]
\varphi(wr,i)&
\end{aligned}\right\}\ \text{par a)}
\]
Ceci dit, $\alpha$ est \unsure{évidemt|évidemment} injectif, son image est formé des autom. \deleted{\unsure{$(u_W,u_R$|$(u_W,u_R)$}} de $\Phi(\Sigma)=(W,R,I,\varphi)$ induisant l'identité sur $I$, i.e. de la forme $(u_W,u_R,\operatorname{id}_I)$, où $u_R$ est un autom affine de $R$, associé à l'autom $u_W$ du grpe $W$ (la \textquotedblleft partie homogène\textquotedblright\ de $u_R$) et satisfaisant
\[
\begin{aligned}
\unsure{(*)|(\alpha)|(\ast\ast)}\quad \deleted{\unsure{*}}\qquad
\varphi(u_Rr,i)&=u_W\,\varphi(r,i)\quad\text{i.e.}\\
\varphi(r,i)&=u_W^{-1}\,\varphi(u_Rr,i)
\end{aligned}
\]
Choisissons en effet $r_0\in R$ [\deleted{\unsure{on aura tout}} \inserted{permettant d'identifier} $R$ à $W$] donc on aura
\[
u_R(wr_0)=u_W(w)ar_0\quad\text{où }\deleted{\unsure{a}}u_R\unsure{(r_0)|wr_0)}=ar_0
\]
et la condition \unsure{$(*)$|$(\alpha)$|$(\ast\ast)$} \deleted{\unsure{a}} s'écrit (pour $r=wr_0\sim w$
% archive sheet number: 3

\scanpage{5}
\[
\begin{aligned}
\varphi(u_W(w)ar_0,i)&=u_W(\varphi(wr_0,i))\qquad\text{ou simplement}\\
\deleted{\unsure{\varphi(u(w)ar_0,i)=u(\varphi(wr_0,i))}}&\\
\Vert\qquad &\qquad\unsure{\leftarrow|\Vert}\\
\deleted{\unsure{\operatorname{int}(u(w))\,\varphi(ar_0,i)=u\varphi}} & u(w\,\varphi(r_0,i)\,w^{-1})\\
u(w)a\,\varphi(r_0,i)a^{-1}u(w)^{-1}&\qquad\Vert\\
&u(w)\,u(\varphi(r_0,i))\,u(w)^{-1}
\end{aligned}
\]
\deleted{donc} ou encore
\[
 a\,\varphi(r_0,i)a^{-1}=u\,\varphi(r_0,i)\qquad\forall\,i\in I
\]
i.e. $u_W=\operatorname{int}(a)$ (car les $\varphi(r_0,i)$, \inserted{$i\in I$,} engendrent $W$

donc
\[
u_R(wr_0)=(awa^{-1})(ar_0)=a(wr_0)=a_R(wr_0)
\]
i.e.
\[
u_R=a_R
\]
i.e.
\[
(u_W,u_R,\operatorname{id}_I)=\alpha(a).
\]

On a \deleted{\unsure{donc}} un diagramme de suites exactes
\diagram{\begin{tikzcd}[ampersand replacement=\&,column sep=small]
\& \& \& \substack{\text{\unsure{gr|gp|gr.} des autom ext}\\\text{de }\Sigma}\arrow[d,equal] \&\\
1\arrow[r] \& W\arrow[r,"\alpha_0"]\arrow[d,Rightarrow,"\mathrm{id}"] \& \operatorname{Aut}(\Sigma)\arrow[r,"\beta_0"]\arrow[d] \& \operatorname{Aut}(I)\arrow[r]\arrow[d,Rightarrow,"\mathrm{id}"] \& 1\\
1\arrow[r,dash] \& W\arrow[r,"\alpha"] \& \operatorname{Aut}(\Phi(\Sigma))\arrow[r,"\beta"] \& \operatorname{Aut}(I) \&
\end{tikzcd}}
et le diagramme des 5 \unsure{mor.} montre alors que la flèche verticale médiane est un iso, cqfd.

\begin{center}\rule{2.5cm}{0.35pt}\end{center}

Les \deleted{\unsure{parties}} \deleted{\unsure{ens}} paraboliques de racines de $\Sigma\subset V$ \inserted{(vectoriel réel de rang \unsure{$l$|1})} correspondent \deleted{\unsure{ent}} \deleted{aux} par \deleted{\unsure{orthogonalité}} \inserted{polarité} aux \uline{facettes} \inserted{\uline{fermées}} du syst. d'hyperplans dans $\check V$ associé au syst. de racines dans $V$.
% archive sheet number: 4

\scanpage{6}
L'inclusion des parties paraboliques (ou encore, des grpes paraboliques correspondants, \unsure{quand|qnd} le syst. de racines $\Sigma$ provient d'un groupe alg) est alors transformée \deleted{par} \inserted{en l'opposée de} l'inclusion des facettes dans $\check V$. C'est cette relation (\uline{opposée} de l'inclusion des paraboliques) qu'on prend sur cet \deleted{\unsure{e}} ens de parties, qui devient dès lors une géométrie de \uline{drapeaux stricts}, correspondant à un \unsure{squelette|complexe} simplicial, dont l'ens des sommets \inserted{(ou simplexes minimaux)} correspond aux \textquotedblleft paraboliques \unsure{maximaux|minimaux}\textquotedblright, et l'ens des simplexes maximaux aux \textquotedblleft Borels\textquotedblright\ ou \textquotedblleft syst simples de racines\textquotedblright\ ou repères $r\in R$. En fait, il s'agit manifestement d'une triangulation de la sphère \inserted{$S^{l-1}$}. L'ens des repères s'identifiera à $R\times\deleted{\unsure{\mathrm{Ord}}}\inserted{\operatorname{Rep}(I)}$, où $\operatorname{Rep}(I)=\operatorname{Bij}([1,l],I)$ est l'ens des ordres totaux sur l'ens $I$ (de \unsure{cart.|cnt.|crt.|card.} $l$).
% archive sheet number: 5

\scanpage{7}
qui est un torseur \inserted{à droite} \unsure{sous|sur} $\unsure{\mathfrak S|\mathscr S}_l$. Les op. $\sigma_0,\ldots,\sigma_l$ sur $R\times\operatorname{Rep}(I)$ s'obtiennent manifestement ainsi : $\sigma_0,\sigma_1,\ldots,\sigma_{l-1}$ sont les opérations $\operatorname{id}_R\times\rho_{\tau_{i+1}}$ ($0\leq i\leq l-2$) où $\tau_{i+1}\in\mathfrak S_l$ est la transposition de $i+1$ et $i+2$ ; \inserted{i.e.}
\[
\sigma_i(r,u:[1,l]\overset{\sim}{-}I)=(r,u\circ\tau_{\unsure{i|i'}})\qquad 0\leq i\leq l-1
\]
quant à $\sigma_{l-1}$, on obtient
\[
\sigma_{l-1}(r,u:[1,l]\xrightarrow{\sim}I)=\{\varphi(r,u(l))\mathbin{.}r,u\}
\]
On trouve bien que \inserted{$\sigma_i^2=1$ ($0\leq i\leq l-1$) et que} $\sigma_i$ $\sigma_j$ commutent si $j\geq i+2$ ($0\leq i,j\leq l-1$) -- le seul cas non \unsure{tt à ft|tout à ft} trivial étant $j=l-1$, $i\leq l-3$,
\[
\begin{aligned}
\sigma_i\sigma_{l-1}(r,u)&=\sigma_i(\varphi(r,u(l)).r,u)=(\varphi(r,u(l)).r,u\circ\tau_{i+1})\\
\sigma_{l-1}\sigma_i(r,u)&=\sigma_{l-1}(r,u\circ\tau_{i+1})
=\varphi(r,\underbrace{(u\circ\tau_{i+1})(l)}_{\substack{\Vert\\u(l)\ \text{si}\\i\leq l-3\ \text{i.e.}\\i+1\leq l-2}}).r,u\circ\tau_{i+1})
\end{aligned}
\]
Ici $(\sigma_i\sigma_{i+1})^3=\operatorname{id}$ -- distinguer d'abord le cas $0\leq i\leq l-2$ -- où celà résulte de $(\tau_{i+1}\tau_{i+2})^3=1$, puis $i=l-1$, où $\sigma_{l-1}\sigma_{l-2}$ est donné par
\[
\begin{aligned}
\sigma_{l-1}\sigma_{l-2}(r,u)&=\sigma_{l-1}(r,u\circ\tau_{l-1})=\varphi(r,(u\circ\tau_{l-1})(l)).r,\unsure{u\circ\tau_{l-1}}\\
\pi(r,u)&=(\underbrace{\overbrace{\varphi(r,u(l-1))}^{w}.r}_{r'},\underbrace{u\circ\tau_{l-1}}_{u'})\\
\pi^2(r,u)&=\inserted{\pi(r',u')=}\deleted{\unsure{\varphi(\varphi(r,u(l-1)).r,\ ,u|\varphi(\varphi(r,u(l-1))r,u)}}\\
&(\underbrace{\varphi(r',\underbrace{u'(l-1)}_{u(l)}).r'}_{\varphi(wr,u(l))\,w\,r,\,u},\underbrace{u'\circ\tau_{l-1}}_{u})
\end{aligned}
\]
% archive sheet number: 6

\scanpage{8}
\[
\begin{aligned}
\pi^2(r,\unsure{\bar u|\pi|u})&=(w\varphi(r,u(l))w^{-1}wr,u)=(w\varphi(r,u(l))r,u)\\
&=(\underbrace{\varphi(r,u(l-1))\varphi(r,u(l))\,r}_{\deleted{\unsure{w'}}\quad r''=w'r},u)\\
\pi^3(r,u)&=(\underbrace{\varphi(r'',u(l-1))r''}_{w'\varphi(r,u(l-1)){w'}^{-1}\underbrace{r''}_{w'r}=w'\varphi(r,u(l-1)).r},u\circ\tau_{l-1})\deleted{\unsure{=\varphi(|\text{cqf}}}
\end{aligned}
\]
\[
\begin{aligned}
&\deleted{\pi^3(r,u)}\\
&\deleted{\pi^{2n}(r,u)=}\\
\pi^3(r,u)&=(\underbrace{\varphi(r,u(l-1))\varphi(r,u(l)).\varphi(r,u(l-1)).r}_{r'''=w''r},\underbrace{u\circ\tau_{l-1}}_{u'})\\
\pi^4(r,u)&=(\underbrace{\varphi(r''',u'(\underset{\deleted{\unsure{l}}}{l-1}))r'''}_{w''\varphi(r,u(l))\deleted{\unsure{w}}{w''}^{-1}w''r},\underbrace{u'\circ\tau_{l-1}}_{u})\\
&=(w''\varphi(r,u(l)\deleted{\unsure{a}})r,u)\\
&=\big[(\varphi(r,u(l-1))\varphi(r,u(l)))^2r,u\big]
\end{aligned}
\]

Soit
\[
\begin{aligned}
\varphi(r,u(l-1))&=a\deleted{\unsure{c}}\\
\varphi(r,u(l))&=b
\end{aligned}
\]
On trouve par récurrence
\[
\left\{\begin{aligned}
\pi^{2n}(r,u)&=(\underbrace{(ab)^nr}_{r_n=w_nr},u)\\
\pi^{2n+1}(r,u)&=(\underbrace{\big((ab)^na\big)r}_{r_n'=w_n'r},\underbrace{u\circ\tau_{l-1}}_{u'})
\end{aligned}\right.
\]
Donc
\[
\pi^{\deleted{\unsure{2n}}i}\deleted{=\mathrm{id}}(r,u)=(r,u)\iff
\left\{\begin{aligned}
i&\equiv0\quad(2)\\[-2pt]
&\text{i.e. }i=2n\\
(ab)^n&=1
\end{aligned}\right.
\]
\margin{i.e. \uuline{$n_{u(l-1),u(l)}\mid n$}

\deleted{i.e. $i\equiv0\ (2n$}

i.e. $\unsure{u|i|n}\equiv0\ (2n_{u(l-1),u(l)})$}
\[
\begin{aligned}
\big[\ \pi^{\unsure{2n+1|2n+2}}(r,u)&=(\underbrace{\varphi(r'_n,u'(l-1)).r'_n}_{\underbrace{w'_n}_{(ab)^na}\,\underbrace{\varphi(r,u(l))}_{b}{w'_n}^{-1}w'_nr},\underbrace{u'\circ\tau_{l-1}}_{u})\\
&=\deleted{\unsure{\varphi(w'_n\varphi(r,u(l))}}\ (\underbrace{(ab)^{n+1}r}_{r_{n+1}=w_{n+1}r},u)\\
\pi^{\unsure{2n+3|2n+2}}(r,u)&=(\underbrace{\varphi(r_{n+1},u(l-1)).r_{n+1}}_{w_{n+1}\varphi(r,u(l-1))w_{n+1}^{-1}w_{n+1}r=(ab)^{n+1}a\,r},u\circ\tau_{l-1})\ \big]
\end{aligned}
\]
% archive sheet number: 7

\scanpage{9}
\[
\pi^i(\unsure{r|-r},u)=(r,u)\quad\forall(r,u)\in R\times\operatorname{Rep}(I)
\iff\left\{\begin{aligned}
&i\equiv0\ (2)\quad\text{i.e. }i=2n\\
&(\varphi(r,u(l-1))\varphi(r,u(l)))^n=1\\
&\hspace{1em}\text{i.e. }\deleted{\unsure{n_{u(l-1),u(l)}\mid n}}\\
&\forall r,u
\end{aligned}\right.
\]
i.e. $n$ multiple du ppcm \inserted{$\unsure{\nu|v}$} des $n_{ij}$ ($i,j\in I$ \inserted{$i\ne j$})

Donc on trouve
\[
\left\{\begin{aligned}
&\text{\uline{Ordre de $\pi$}}=\sigma_{l-2}\sigma_{l-1}\qquad=2\nu\\
&\text{où }\nu\text{ est le ppcm des }n_{ij}\quad(i,j\in I,i\ne j).
\end{aligned}\right.
\]
% archive sheet number: 8

\scanpage{11}
\uline{Théorème} (à prouver)

\textcircled{1} Toute carte cellulaire \inserted{$\ne\varnothing$} connexe et simplement connexe quasi régulière (i.e. dont tous les sommets comb. et toutes les faces ont même ordre) est régulière, et elle est complète (i.e. les sommets comb. d'ordre fini sont effectifs) sauf dans le cas d'une carte isomorphe à $(\mathbb R^2\supset\mathbb R)$ (dont la \unsure{complétion est|complétion|complétée est} la 2sphère munie d'un grand-cercle ponctué\ldots).

\textcircled{2} Une \unsure{Carte|carte} cellulaire \inserted{$\ne\varnothing$} conn. simpl. connexe régulière et complète (appelons \deleted{les} \inserted{une telle carte} \textquotedblleft cartes standart\textquotedblright) est déterminée, à isomorphisme près, par \inserted{le couple $(p,q)$} \inserted{où $p$ est} l'ordre des sommets comb., $q$ l'ordre des faces ($p,q\in\overline{\mathbb N^*}$, entiers $1\leq p,q\leq+\infty$). Le couple $(p,q)$ peut prendre toutes les valeurs \deleted{sauf} $\in\overline{\mathbb N^*}\times\overline{\mathbb N^*}$ \uline{sauf} $(1,\unsure{q|p})$ \inserted{et $(p,1)$ si $p\ne2$}. Appelons $(C_{p,q},r_{p,q})$ la carte standard épinglée (i.e. munie d'un biarc-repère $r_{p,q}$) correspondant à $(p,q)$, déterminée à \deleted{homéomorphisme} \inserted{isom. unique} près dans la catégorie isotopique des cartes ]

\textcircled{3}
\[
\operatorname{Aut}(C_{p,q})\xleftarrow{\sim}
\underset{\text{\unsure{groupe cartographique|gpe cartographique}}}{{\unsure{\boldsymbol{\Gamma}|\Gamma}}^{0}}/
(\rho_s^p,\rho_f^q)\overset{\text{déf}}{=}\boldsymbol{\Gamma}_{p,q}
\]
\deleted{\textcircled{\unsure{3|4}}}

\textcircled{4} \uline{Corollaire} Soit $\unsure{\mathscr C|C|\mathcal C}=(X\supset K\supset S)$ une carte cellulaire complète \inserted{connexe} quasi-régulière de type $(p,q)$. Alors le revêtement universel \inserted{$\unsure{\overline C|\widetilde C|\widetilde{\mathcal C}}$} est \deleted{\unsure{rég}} une carte standart, le groupe fondamen%
% archive sheet number: 10
\scanpage{12}%
tal de $X$ s'identifie à un sous-groupe de $\deleted{\unsure{\Gamma}}\Gamma_{p,q}$. Si $\unsure{\mathscr C|C|\mathcal C}$ est régulière \inserted{épinglée}, $G$ le grpe de ses automorphismes, alors le grpe fondamental mixte de $(X,G)$ s'identifie à $\Gamma_{p,q}$

\textcircled{5} [Appelons \deleted{\unsure{plan métrique espace|plan métrique}} \textquotedblleft surface \inserted{(riemannienne)} standart\textquotedblright\ une surface \deleted{\unsure{diff}} \inserted{$E$} analytique réelle munie d'une métrique \inserted{riemannienne $>0$} telle que

a) $E$ est $\ne\varnothing$, connexe, simplement connexe

b) Le groupe des isométries de $E$ est transitif sur $E$.

\deleted{\unsure{c)}} On sait alors que, à un facteur constant près sur la métrique, $E$ est isomorphe à un des trois espaces :

1\textsuperscript{o}) La sphère $S^2$ comme espace homogène de $\mathrm{SO}(3,\mathbb R)$
% slanted
\margin{espace sphérique, déplacements sphériques}

2\textsuperscript{o}) Le plan \deleted{euclidien} \inserted{complexe} \inserted{avec $\unsure{\{z\bar z|qz\bar z|\{z,\bar z\}}$} comme espace homogène sous le groupe $\unsure{\mathbb U|U|\mathcal U}.\unsure{C|\mathbb C}$ (1/2 direct) le groupe des déplacements euclidiens
% slanted
\margin{plan (ou espace) euclidien,}

3\textsuperscript{o}) Le demi-plan de Poincaré, \deleted{\unsure{comme|comme espace}} espace homogène de $\mathrm{SL}(2,\mathbb R)/\{\pm1\}$.
% slanted
\margin{plan hyperbolique, dépl. hyperboliques}

Les groupes écrits sont les \uline{comp. connexes} du groupe des isométries -- ou \unsure{encore|ens|mieux}, \inserted{\unsure{formés|formé|fermées}} des isométries qui \uline{conservent} \uline{l'orientation} -- ou encore \unsure{le|la|les}
% archive sheet number: 11

\scanpage{13}
structure complexe associée à l'orientation (dans le cas où une telle orientation est choisie). Noter que dans les cas 1\textsuperscript{o} et 3\textsuperscript{o}, où les similitudes sont des isométries, on peut \uline{normaliser} la métrique -- on parlera alors de \textquotedblleft surface standard normalisée\textquotedblright\ -- déterminée, pour chacun des \deleted{\unsure{deux|trois}} trois types, à \uline{isom} près.

Une carte cellulaire sur une surface riemannienne \inserted{$X$} est dite \uline{géodésique} si \deleted{elle} ses arêtes ouvertes sont des \inserted{arcs de} géodésiques ; elle est dite \uline{métriquement régulière} si elle est géodésique, et si le groupe des automorphismes \uline{isométriques} de la carte est transitif sur l'ens des répères (donc simplement transitif sur l'ens des répères) On dit que c'est une \uline{carte géodésique \unsure{standart|standard}} si elle est métriquement régulière, \inserted{complète (cf. 1\textsuperscript{o}),} et si \deleted{\unsure{la}} $X$ \deleted{\unsure{y}} est une surface riemannienne \deleted{\unsure{normalisée}} \inserted{standart} \unsure{?} ]

\deleted{On dit qu'elle est} \deleted{\unsure{de si|de type}} \inserted{Ceci posé, on doit avoir ceci} ]

a) Pour toute carte standart \inserted{$(X\supset K\supset S)$}, il existe \deleted{\unsure{sur|une}} une carte géodésique standard $(X_0,K_0,S_0)$ qui lui est (topologiquement) isomorphe.

\deleted{\unsure{Bi}} b) Pour deux cartes géodésiques standard, dans tte classe d'isotopie d'isomorphismes
% archive sheet number: 12

\scanpage{14}
topologiques il y a un isomorphisme topologique unique qui soit une \deleted{\unsure{isomé}} \inserted{similitude}. (donc une isométrie dans le cas où les surfaces sous-jacentes sont \deleted{\unsure{hémiques|hériques|hyperboliques}} sphériques normalisées ou hyperboliques normalisées)

\inserted{c) Le groupe \unsure{des} autom. métriques d'une carte géod. st. \unsure{$(X,\text{etc.})$|$(X,K,S)$} \inserted{\unsure{métrique} $(X\supset K\supset S,g)$} opère \uline{proprement} sur $X$ -- il \deleted{\unsure{est}} c'est un ss-grpe \uline{discret} de $\operatorname{Aut}(X,g)$}

\uline{NB} On dit que la carte géodésique standard est \uline{normalisée} si a) \unsure{Dans|dans} le cas sphérique ou hyperbolique, $(X,g)$ est normalisé b) dans le cas euclidien -- où, cf plus bas, la carte est \unsure{loc.|loc. f.} finie \deleted{\unsure{sommets et faces,}} donc les longueurs des arêtes sont toutes égales -- on demande que la longueur des arêtes soit égale à $1$ (on pourrait \deleted{\unsure{aussi}} \inserted{à la place} exiger que l'aire des faces soit égale à $1\ldots$). On peut donc \deleted{dire que} \inserted{\uline{Corollaire}} Le foncteur \textquotedblleft oubli de la métrique\textquotedblright\ de la catégorie des cartes \deleted{cellulaires} \inserted{géodésiques} \deleted{\unsure{g}} standard normalisées (avec comme morphismes les isométries qui respectent la carte) dans la catégorie des cartes standard, est une équivalence de catégories.

\uline{Re NB} Carte géodésique standard normalisée \uline{épinglée} : \unsure{Les|Les deux} surfaces ambiantes s'\unsure{identifie-|identifie.|identifient}%
% archive sheet number: 13

\scanpage{15}
canoniquement aux surfaces riemanniennes standart \uline{types} ${\unsure{\mathbb S|S}}^2,\mathbb R^2=\mathbb C,\mathbb H$ avec le \textquotedblleft répère\textquotedblright\ différentiel habituel \ldots. On utilise \unsure{comme|encore} notation $C_{p,q}$

\deleted{\unsure{6}} On distingue \deleted{types} des cartes standard de type sphérique, euclidien ou hyperbolique.

\textcircled{6} \textcircled{a} \deleted{\unsure{p,q}} $C_{p,q}$ est sphérique \deleted{\unsure{sss}} exactement dans les cas suivants : \deleted{\unsure{$p=2,q$|$p=q$}} $p=2$, (\deleted{\unsure{$q$ quelconque $\geq1$}}\inserted{$1\leq q<+\infty$}) ; $q=2$ (\deleted{\unsure{$p$ quelconque $\geq1$}}\inserted{$1\leq p<+\infty$}), \deleted{$p=q=3,\ p=$}
\diagram{\(\begin{array}{ccccc}
(3,3)&(3,4)&(4,3)&(3,5)&(5,3)\\
\text{tétraèdre}&\text{cube}&\text{octaèdre}&\text{dodécaèdre}&\text{icosaèdre}
\end{array}\)}
-- ou encore sss $C_{p,q}$ \uline{finie} ou encore sss $\Gamma_{p,q}$ fini.

\textcircled{b} $C_{p,q}$ est euclidien exactement dans les cas suivants : \inserted{$(2,\infty)$, \deleted{\unsure{$(\infty,$|$(\infty,2)$}} $(\infty,2)$}
\diagram{\(\begin{array}{ccc}
(4,4),&(3,6),&(6,3)\\
\text{pavage carré}&\text{pavage}&\text{pavage}\\
&\text{hexagonal}&\text{triangulaire}
\end{array}\)}
\deleted{\unsure{$(2,\infty)$}} -- ou encore sss $\Gamma_{pq}$ est infini et \uline{résoluble}.

\textcircled{c} $C_{pq}$ est hyperbolique dans les autres cas i.e. \deleted{\unsure{$(p,q)>(3,3)$}} \deleted{\unsure{$(p>3,\ q\geq3)$|$(p\geq3,\ q\geq3)$}} \deleted{\unsure{ou $(p=3$,}}
\[
\text{\textcircled{$c_1$}}\quad
\left\{\begin{aligned}
&\text{ou }\deleted{\unsure{(p\geq3,q\geq3,\ (p,q)\ne(4,4))\ \text{ou}}}\qquad(p,q)>(4,4)\\
&\text{ou }\quad p=3,\ q\geq7\\
&\text{ou }\quad q=3,\ p\geq7
\end{aligned}\right.
\]
\margin{ordre produit sur $\unsure{\overline{\mathbb N}^{\,*}\times\overline{\mathbb N}^{\,*}|\mathbb N^*\times\overline{\mathbb N}^{\,*}}$}
\[
\begin{aligned}
\text{\textcircled{$c'_1$}}\ &\text{ou encore }\quad p+q\geq10\ \text{ou }(p,q)=(4,5)\ \text{ou }(p,q)=(5,4)\\
\unsure{(c''_1)|(c'_1)}\ &\text{ou encore }\quad p+q\geq9\ \text{et }(p,q)\ne(3,6),(6,3)\quad\deleted{\unsure{\text{grpe}}}\\
(c_2)\ &\text{ou encore }\quad\Gamma_{pq}\ \text{infini non résoluble}\quad\deleted{\unsure{(?)}}
\end{aligned}
\]
$(c'_2)$ ou encore $\Gamma_{pq}$ contient ss-groupe libre à 2 générateurs (donc contient \unsure{un} groupe libre à $n$ générateurs $\forall n\geq0$)

\margin{$\checkmark$}
\textcircled{7} Les cartes \inserted{quasi-}régulières \inserted{$(X\supset K\supset S)$} épinglées \inserted{de type $(p,q)$} s'obtiennent (à isom canonique près dans la cat. isotopique des cartes) à partir des sous-groupes $\pi\subset\Gamma_{p,q}$
% archive sheet number: 14

\scanpage{16}
\uline{opérant} \deleted{sans points fixes} \inserted{\uline{librement}} sur $X_{pq}$, en \unsure{prenant le|passant au} quotient de $C_{pq}$ par $\pi$ (qui s'identifie au groupe fondamental de $X\ldots$) Elles sont donc munies (canoniquement modulo isotopie de la carte) d'une métrique riemannienne (la métrique quotient) pour laquelle la carte est \uline{géodésique}. \deleted{\unsure{et}} Elle est régulière sss le sous-groupe $\pi$ de $\Gamma_{pq}$ est invariant.
% slanted
\margin{\uuline{attention} au cas où $p$ ou $q$ est infini ! \unsure{(passant|passant} au quotient, il peut y avoir création de nouveaux sommets comb. d'ordre fini\ldots}
% slanted
\margin{Notion de carte \uline{métriquement} quasi-régulière (\unsure{complète|complétée} \deleted{\unsure{et quasi-rég}}, i.e. dont le rev. universel est carte géod. standart) \unsure{Donner|Donnes} encore th. d'équiv. de cat. (\uline{sans} \unsure{repèr|répèr} \ldots)}

\uline{NB} Il y a donc lieu, pour \inserted{tt} $(p,q)$ donné, pour déterminer les $\pi$ possibles, d'étudier \deleted{\unsure{quels}} l'op. de $\Gamma_{pq}$ sur $X_{pq}$ -- et plus particulièrement, de trouver les \deleted{pts de $X_{pq}$} \deleted{\unsure{groupes de}} stabilisateurs dans $\Gamma_{pq}$ des pts de $X_{pq}$.

\deleted{\unsure{Corollaire}} \uline{Re NB} \unsure{Si|Il} \deleted{\unsure{on a|y a}} la carte \inserted{qu. rég.} est orientée \inserted{(p.ex. \uline{\unsure{orient.bl|orient.ble|orientable}} et épinglée)}, \deleted{\unsure{elle}} son représentant métriquement quasi-régulier se trouve muni canoniquement d'une structure \uline{complexe} -- qui, dans le cas compact \inserted{(i.e. carte \uline{finie})} \inserted{(i.e. $p,q$ finis et $\pi$ d'\uline{indice} \uline{fini})} en fait donc une \uline{courbe algébrique} projective \unsure{de ma suite|de ma récolte|de marsault|de ma seule} ! \unsure{Sa} métrique est celle de Bergman.
% archive sheet number: 15

